%% kapa: exam=1 round=37 year=2026 period=1 form=A subject=economics subject_ko=경제학원론 seq=21 number=61
%% source: Q-Net 감정평가사 기출문제 (cst00309) artlSeq=5258924 file=raw/1008/5258924/2026년 제37회 감정평가사 1차 1교시.pdf
%% conversion: kordoc 4.24.51 -> tools/gen1.py (auto); review=proofread-sonnet; flags=none
\begin{question}{61}
국민소득 지표에 관한 설명으로 옳은 것을 모두 고른 것은?
\begin{framedtext}
ㄱ. 외국인의 한국 내 생산활동은 한국의 GDP(국내총생산)에 포함되지 않는다. \par
ㄴ. 폐쇄경제에서는 실질 GDP와 실질 GDI(국내총소득)가 같다. \par
ㄷ. GNP(국민총생산)는 GDP에서 국외 순수취 요소소득을 뺀 수치이다. \par
ㄹ. 정부의 출산장려금은 GDP에 포함되지 않는다. \par
\end{framedtext}
\begin{choices}
  \item ㄱ, ㄴ
  \item ㄱ, ㄷ
  \item ㄴ, ㄷ
  \item ㄴ, ㄹ
  \item ㄷ, ㄹ
\end{choices}
\end{question}
\answer{4}
